\begin{trapbox}
	\begin{description}[style=nextline, leftmargin=2.8em, labelsep=0.5em, itemsep=0.35em]
		\item[\textbf{\textcolor{CTrap}{易错点一（忽略定义域）}}] 错误认为 $n$ 可以为任意实数，将公式应用于 $n$ 为非正整数情形。
		\item[\textbf{\textcolor{CTrap}{正确理解（明确条件）：}}] 公式只对正整数 $n$ 成立，因为阶乘定义在非负整数上。
		\item[\textbf{\textcolor{CTrap}{错因分析（定义域不清）：}}] 学生常忽视阶乘的定义域，将 $n!$ 误当作一般函数。
	\end{description}

	\medskip
	\begin{description}[style=nextline, leftmargin=2.8em, labelsep=0.5em, itemsep=0.35em]
		\item[\textbf{\textcolor{CTrap}{易错点二（抵消项错误）}}] 在求和 $S_n = \sum_{k=1}^{n} k\cdot k!$ 时，错误认为结果等于 $n! - 1$ 或 $(n+1)!$。
		\item[\textbf{\textcolor{CTrap}{正确理解（首尾分析）：}}] 正确应为 $(n+1)! - 1! = (n+1)! - 1$，消去中间项后只剩末项减首项。
		\item[\textbf{\textcolor{CTrap}{错因分析（裂项展开不全）：}}] 未正确写出第一项和最后一项，展开时丢项或符号错误。
	\end{description}

	\medskip
	\begin{description}[style=nextline, leftmargin=2.8em, labelsep=0.5em, itemsep=0.35em]
		\item[\textbf{\textcolor{CTrap}{易错点三（裂项方向倒置）}}] 错误认为 $n\cdot n! = n! - (n-1)!$。
		\item[\textbf{\textcolor{CTrap}{正确理解（定义推导）：}}] 由 $(n+1)! = (n+1)n!$ 得 $(n+1)! - n! = n\cdot n!$。
		\item[\textbf{\textcolor{CTrap}{错因分析（定义混淆）：}}] 未掌握阶乘的递推关系，将 $(n+1)!$ 错误写为 $n!(n-1)!$ 等。
	\end{description}
\end{trapbox}
